% ============================================================
\section{Methods}
% ============================================================

The audit establishes what the rules are; it cannot show that the partial ones
matter, and it cannot show that the two coordinates come apart under
intervention rather than only on paper. We therefore built the smallest system
in which both coordinates can be manipulated independently, and ran the three
arms in it. Code, configurations, seeds and the analysis pipeline are
available in the repository described below.\footnote{Repository: see Data
Availability. Every result file carries a
manifest with the full run specifications, a hash of the source that produced
it, and the environment versions.}

\subsection{Substrate}

A vector-genome ecology of digital organisms, written in \textsc{numpy}. We
deliberately did not build an instruction-set interpreter in the Tierra/Avida
lineage. The hypotheses require heredity, variation, differential reproduction,
competition for a limited resource, and the ability to label lineages and
condition allocations on the label. They do not require programs that copy
themselves, and an interpreter would have added a validation burden that buys
nothing here.

Each organism carries an energy reserve, a lineage label $\ell$, and four
genes, each with a declared role and a quantitative prediction attached to it:
acquisition affinity $a$, whose share of the resource is proportional to $a$
and whose maintenance cost is $c_0 + k_c a^2$; the fraction $\varphi$ of the
parent's reserve transferred to offspring; a modifier of the organism's own
mutation rate; and a selectively neutral marker.

Two design constraints were treated as binding. First, \emph{every gene must do
something}: a gene with no effect turns everything measured on it into drift.
Second, \emph{every gene must face a trade-off}. Without the quadratic cost
term, $a$ increases without bound; without a cost of holding reserves,
$\varphi$ runs to its ceiling because a parent can wait indefinitely at no
risk. Both failures occurred in development and both were fixed by adding the
missing cost, not by bounding the gene.

Population density is self-regulating rather than imposed: the total resource
per step is fixed, so per-capita income falls as the population grows, and
mortality is condition-dependent, falling with the reserve. There is no
arbitrary culling rule and no carrying-capacity parameter.

\subsection{Substrate validation}

None of what follows is interpretable if the world does not behave like an
evolving population. We therefore validated it against four quantitative
predictions external to the code that evaluates them, and treated this as
blocking. The strongest is specific to this model: the affinity gene has an
interior evolutionarily stable value, obtained by setting the invasion gradient
to zero in a monomorphic population of size $N$,
\begin{equation}
a^{*} \;=\; \sqrt{\frac{R\,(N-1)}{2\,k_c\,N^{2}}} ,
\label{eq:ess}
\end{equation}
which is a sharp prediction that varies with two parameters. The others are the
breeder's equation and drift. A fourth, mutation--drift balance, is a
scale-consistency check on the implementation rather than an independent test
of drift theory, and we report it as such. To these four we add a robustness
check, the insensitivity of the qualitative pattern to arbitrary parameters:
five checks in all. Results are in \S6.1.

Two of these tests are easy to run wrongly, and we report how because the wrong
version passes inspection. Drift must not be measured with a marker that starts
variable alongside a gene under selection: what is then observed is genetic
draft, it decays far faster and erratically, and it does not resemble
$(1-1/2N)^g$. We freeze the physiological genes for that test. And the selection
differential must be measured against the set of candidates the truncation acts
on, not against the whole population; against the whole population $S$ is
diluted by a demographic factor and $R/S$ departs from unity with nothing
wrong.

\subsection{The manipulated rule and the four conditions}

The rule we manipulate allocates energy between organisms. Collection is
identical in all four conditions: every organism above a threshold contributes
a fixed amount with fixed probability. Allocation sends each contribution to
the neediest candidate in a random sample. The conditions are the four ways of
varying two things and nothing else:

\begin{description}[leftmargin=*,style=nextline,itemsep=1pt]
\item[$C$:] whether the transfer is executed by the organisms themselves
($\mathrm{I}$) or by a central allocator ($\mathrm{E}$).
\item[$A$:] whether the candidate set is any living organism ($A=0$) or is
conditioned on the label ($A=1$)---the focal label for the central allocator,
the donor's own label for the organisms.
\end{description}

Volume, timing, criterion, sample size and topology are identical across the
four cells by construction, and the random number stream is consumed in the
same order, which is what makes the regression tests below exact rather than
approximate.

\subsection{The two interventions}

\paragraph{Removal of the external operator}
The support supplied from outside $D$ is withheld. Where $C=\mathrm{E}$ the
transfer stops. Where $C=\mathrm{I}$ there is no external machinery to withhold:
the mechanism is instantiated inside the focal system, and removing it would be
amputating part of $D$ rather than withdrawing support from it. That asymmetry
is not an artefact of the implementation---it is what the control coordinate
asserts, and the intervention is simply the coordinate made operational.

\paragraph{Relabelling}
Once the population is stable, living organisms are cloned: state copied
exactly, label changed. This is the test of \S3.2 run inside a controlled
system, and it settles by construction what the test elsewhere has to
establish by assumption---that the two configurations differ in the label and
in nothing else.

The observable is what the rule \emph{allocates}, per organism-step, within a
window after the injection. It is not long-run survival of the injected
lineage, and the difference is not cosmetic. Under an exchangeable rule the
injected lineage is a neutral marker, so its frequency drifts to absorption
with fixation probability equal to its initial frequency; averaging the final
population over seeds then returns the initial fraction and looks exactly like
persistence while being its opposite. An earlier version of this experiment
reported that average as a result. The per-seed contrast behind it was one
survival in ten against zero in ten, which Fisher's exact test does not
distinguish ($p=1.0$).

\subsection{The three arms}

Three manipulations were run, each with arms B, R and S as defined in \S3.4,
with $\Phi$ the autonomous reproductive cycles completed by the focal lineage
within a fixed horizon, and with the size of a competing lineage varied as a
dose.

The first is a negative control. The audit says the transfer rule under
$\mathrm{E/p}$ is recipient-partial; the pre-registered prediction is
nonetheless $\rho=\texttt{NA}$, because baseline dependence is not positive. A
method that returned a small $\rho$ there would be manufacturing a measurement
where there is no dependence to apportion. The second supplies a reproductive
resource, partially in arm B and under an open, dynamically attainable
enrolment condition in arm S. The third suppresses competitors: arm B bars the
non-focal lineage from reproducing, and arm S bars the same fraction by a
symmetric condition on state. The third is where $\rho$ should be near one,
because the benefit in arm B is \emph{being the exempt one}, and an open rule
cannot reproduce it---it would bar the focal lineage too.

Building an exchangeable provider is itself a design act, and a rule that is
formally open can be practically closed. The enrolment condition is therefore
dynamically attainable in the sense of \S2.2---any organism with enough reserve
can pay it---and the relabelling test is available as a check: if no foreign
lineage could ever qualify, the rule was partial all along. Allocating
computational resources through open markets is long established outside this
field \citep{sutherland1968,miller1988,waldspurger1992}; we borrow the idea of
an access rule indifferent to who is asking rather than proposing it. What we
have not found is its use as a controlled substitution arm in an evolutionary
system.

\paragraph{Matching}
Arm S was matched to arm B on emitted volume for the resource manipulation and
on suppressed fraction for the suppression manipulation, calibrated in pilot
runs whose seeds are disjoint from the confirmatory ones and which are
discarded. Timing, latency and topology are identical by construction. The
acquisition cost the recipient pays in arm S is reported separately and never
absorbed into $\Phi$. The tolerance, the guard $\epsilon$, the number of seeds
and the primary definition of $\Phi$ were fixed before the confirmatory runs.

\paragraph{What we did not run}
One manipulation suggested by \S4.5 is absent: making execution itself a
service, so that any entity may step any other on a rule open to all. It is the
arm closest to the on-chain systems, and it remains the most valuable thing to
run next. Its central claim, however---that two populations equally dependent
on external supply can differ in \emph{for whom} that supply exists---is
exactly what the $\mathrm{E/x}$ versus $\mathrm{E/p}$ contrast tests here,
under a rule we could hold identical in every other respect. We took the
general form first.
